% GENERATED FILE. Edit canonical lessons, then run npm run generate:books.
% canonical-source-hash: fnv1a64:01a6cae176241d54
% canonical-lessons: ES-C682-farola, ES-C682-papelera, ES-C682-contenedor, ES-C682-bordillo, ES-C682-peatonal

\chapter{Lo que hay en la calle}
\label{ch:es-lo-que-hay-en-la-calle}
\hlchaptermodality{682}

\begin{chapteropening}
\textbf{By the end of this chapter:} \emph{I can talk about what is in the street: lamp posts, bins, kerbs and pedestrian streets.}

Complete the last lesson of chapter 682: I can talk about what is in the street: lamp posts, bins, kerbs and pedestrian streets.
\end{chapteropening}

\section[la farola]{la farola — the street lamp, the lamp post}

\label{lesson:ES-C682-farola}



\begin{quote}
Before the new one: say the Spanish for holiday, then the Spanish for to serve.
\end{quote}

\subsection*{You'll want to know: la farola}
\textbf{la farola} — \textquotedblleft{}the street lamp, the lamp post\textquotedblright{}.

\begin{quote}
Las farolas se encienden a las nueve.
\end{quote}

\begin{quote}
\emph{The street lamps come on at nine.}
\end{quote}

A Spain word for a tall street lamp.

\begin{cousinweb}
\emph{farola} is built on \emph{farol} \textquotedblleft{}lantern\textquotedblright{}. The word is usually linked to \emph{faro} \textquotedblleft{}lighthouse\textquotedblright{}, from Latin \emph{pharus}, though the details are not settled.
\end{cousinweb}

\begin{grammarlens}[title={What you've built}]
One more word for telling, explaining and describing at B1.
\end{grammarlens}

\subsection*{Your turn}
\begin{itemize}
\raggedright
  \item \emph{Say it:} \emph{la farola}
  \item \emph{Say it:} \emph{la farola}, once more
  \item \emph{Say it:} a sentence of your own with \emph{la farola}
  \item \emph{Recall:} say the Spanish for holiday, then the Spanish for to serve, then say \emph{la farola} again
\end{itemize}

\begin{tcolorbox}[breakable,colback=teal!4,colframe=teal!35!black,title={Before you move on}]
What does la farola mean? (\textquotedblleft{}The street lamp, the lamp post\textquotedblright{}.) Say it once more.
\end{tcolorbox}
\section[la papelera]{la papelera — the litter bin, the wastepaper basket}

\label{lesson:ES-C682-papelera}



\begin{quote}
Before the new one: say the Spanish for to serve, then the Spanish for the street lamp.
\end{quote}

\subsection*{You'll want to know: la papelera}
\textbf{la papelera} — \textquotedblleft{}the litter bin, the wastepaper basket\textquotedblright{}.

\begin{quote}
Tira el papel a la papelera.
\end{quote}

\begin{quote}
\emph{Throw the paper in the bin.}
\end{quote}

For paper and small rubbish; \emph{la basura} is the rubbish.

\begin{cousinweb}
\emph{papelera} is built on \emph{papel} \textquotedblleft{}paper\textquotedblright{}, which came from Catalan \emph{paper}, from Latin \emph{papyrus}, Greek \emph{pápyros} (EtymDB).
\end{cousinweb}

\begin{grammarlens}[title={What you've built}]
One more word for telling, explaining and describing at B1.
\end{grammarlens}

\subsection*{Your turn}
\begin{itemize}
\raggedright
  \item \emph{Say it:} \emph{la papelera}
  \item \emph{Say it:} \emph{la papelera}, once more
  \item \emph{Say it:} a sentence of your own with \emph{la papelera}
  \item \emph{Recall:} say the Spanish for to serve, then the Spanish for the street lamp, then say \emph{la papelera} again
\end{itemize}

\begin{tcolorbox}[breakable,colback=teal!4,colframe=teal!35!black,title={Before you move on}]
What does la papelera mean? (\textquotedblleft{}The litter bin, the wastepaper basket\textquotedblright{}.) Say it once more.
\end{tcolorbox}
\section[el contenedor]{el contenedor — the large bin; the container}

\label{lesson:ES-C682-contenedor}



\begin{quote}
Before the new one: say the Spanish for the street lamp, then the Spanish for the litter bin.
\end{quote}

\subsection*{You'll want to know: el contenedor}
\textbf{el contenedor} — \textquotedblleft{}the large bin; the container\textquotedblright{}.

\begin{quote}
El vidrio va al contenedor verde.
\end{quote}

\begin{quote}
\emph{Glass goes in the green bin.}
\end{quote}

\emph{El contenedor de reciclaje} — the recycling bin.

\begin{cousinweb}
\emph{contenedor} is built on \emph{contener} \textquotedblleft{}to contain\textquotedblright{}, which comes from Latin \emph{continēre} \textquotedblleft{}to hold together, to contain\textquotedblright{} (EtymDB), built on \emph{tenēre} \textquotedblleft{}to hold\textquotedblright{}.
\end{cousinweb}

\begin{grammarlens}[title={What you've built}]
One more word for telling, explaining and describing at B1.
\end{grammarlens}

\subsection*{Your turn}
\begin{itemize}
\raggedright
  \item \emph{Say it:} \emph{el contenedor}
  \item \emph{Say it:} \emph{el contenedor}, once more
  \item \emph{Say it:} a sentence of your own with \emph{el contenedor}
  \item \emph{Recall:} say the Spanish for the street lamp, then the Spanish for the litter bin, then say \emph{el contenedor} again
\end{itemize}

\begin{tcolorbox}[breakable,colback=teal!4,colframe=teal!35!black,title={Before you move on}]
What does el contenedor mean? (\textquotedblleft{}The large bin; the container\textquotedblright{}.) Say it once more.
\end{tcolorbox}
\section[el bordillo]{el bordillo — the kerb, the curb}

\label{lesson:ES-C682-bordillo}



\begin{quote}
Before the new one: say the Spanish for the litter bin, then the Spanish for the large bin.
\end{quote}

\subsection*{You'll want to know: el bordillo}
\textbf{el bordillo} — \textquotedblleft{}the kerb, the curb\textquotedblright{}.

\begin{quote}
El autobús para junto al bordillo.
\end{quote}

\begin{quote}
\emph{The bus stops next to the kerb.}
\end{quote}

A Spain word; elsewhere \emph{el borde de la acera} or \emph{el cordón}.

\begin{cousinweb}
\emph{bordillo} is built on \emph{borde} \textquotedblleft{}edge\textquotedblright{}, which comes from French \emph{bord} (EtymDB), a word of Germanic origin.
\end{cousinweb}

\begin{grammarlens}[title={What you've built}]
One more word for telling, explaining and describing at B1.
\end{grammarlens}

\subsection*{Your turn}
\begin{itemize}
\raggedright
  \item \emph{Say it:} \emph{el bordillo}
  \item \emph{Say it:} \emph{el bordillo}, once more
  \item \emph{Say it:} a sentence of your own with \emph{el bordillo}
  \item \emph{Recall:} say the Spanish for the litter bin, then the Spanish for the large bin, then say \emph{el bordillo} again
\end{itemize}

\begin{tcolorbox}[breakable,colback=teal!4,colframe=teal!35!black,title={Before you move on}]
What does el bordillo mean? (\textquotedblleft{}The kerb, the curb\textquotedblright{}.) Say it once more.
\end{tcolorbox}
\section[peatonal]{peatonal — pedestrian (as an adjective)}

\label{lesson:ES-C682-peatonal}



\begin{quote}
Before the new one: say the Spanish for the large bin, then the Spanish for the kerb.
\end{quote}

\subsection*{You'll want to know: peatonal}
\textbf{peatonal} — \textquotedblleft{}pedestrian (as an adjective)\textquotedblright{}.

\begin{quote}
Esta calle es peatonal.
\end{quote}

\begin{quote}
\emph{This is a pedestrian street.}
\end{quote}

\emph{Calle peatonal} or \emph{zona peatonal} — no cars.

\begin{cousinweb}
\emph{peatonal} is built on \emph{peatón} \textquotedblleft{}pedestrian\textquotedblright{}, which came from French \emph{piéton}, built on \emph{pied} \textquotedblleft{}foot\textquotedblright{}, from Latin \emph{pēs}.
\end{cousinweb}

\begin{grammarlens}[title={What you've built}]
One more word for telling, explaining and describing at B1.
\end{grammarlens}

\subsection*{Your turn}
\begin{itemize}
\raggedright
  \item \emph{Say it:} \emph{peatonal}
  \item \emph{Say it:} \emph{peatonal}, once more
  \item \emph{Say it:} a sentence of your own with \emph{peatonal}
  \item \emph{Recall:} say the Spanish for the large bin, then the Spanish for the kerb, then say \emph{peatonal} again
\end{itemize}

\begin{tcolorbox}[breakable,colback=teal!4,colframe=teal!35!black,title={Before you move on}]
What does peatonal mean? (\textquotedblleft{}Pedestrian (as an adjective)\textquotedblright{}.) Say it once more.
\end{tcolorbox}
